% !TeX root = ../main.tex

\xchapter{燃料中心温度与裂变气体释放模型的棒级验证及适用性分析}{Rod-Level Validation and Applicability Assessment of Fuel Centerline Temperature and Fission Gas Release Models}

\xsection{引言}{Introduction}

第二章至第四章分别对燃料与包壳材料物性、元素及孔隙迁移行为、裂变气体模型、间隙导热模型以及非线性力学本构模型进行了建立与单项验证。这些验证主要面向材料关系、局部本构积分和子模型接口，能够说明各物理模型在独立或局部耦合条件下的实现正确性。然而，在实际燃料棒服役过程中，燃料温度场、裂变气体释放、间隙气体组成、气隙导热和燃料棒内压之间存在显著的相互反馈，仅依靠单项模型验证尚不足以评价程序在棒级多物理场耦合计算中的整体预测能力。因此，本章进一步基于公开基准数据和补充实验数据，开展燃料棒尺度的综合验证与适用性分析。

本章重点围绕燃料中心温度和裂变气体释放两个关键响应量展开验证。燃料中心温度直接反映燃料热导率、间隙传热和边界换热模型的综合作用；裂变气体释放则同时受温度历史、燃耗、裂变率、晶粒尺度气体迁移和晶界释放过程控制，并进一步通过改变气隙气体组成影响间隙热导率。根据 FRAPCON-4.0 验证报告，燃料性能程序常用的实验验证量主要包括燃料温度、裂变气体释放、气腔体积、包壳腐蚀和包壳环向应变等。结合本文研究重点及前文模型验证基础，本章将验证对象限定为 MOX 燃料棒的中心温度和裂变气体释放行为，并在此基础上分析二者与间隙导热之间的耦合反馈机制。

针对裂变气体行为，第四章中已实现工程化简化机理模型 ForMas 和高保真机理模型 SCIANTIX 两类模型。二者在计算代价、输入参数需求、细观机制描述能力和长时程棒级计算适用性方面存在差异，因此有必要在相同棒级验证框架下进行对比分析。本章首先介绍验证数据来源、计算模型设置、边界条件处理和误差评价方法；随后开展燃料中心温度验证，并分别给出基于 ForMas 和 SCIANTIX 的裂变气体释放验证结果；在此基础上比较两类模型的预测精度、计算效率和适用范围；最后分析温度、裂变气体释放和间隙导热之间的耦合反馈关系，为后续燃料棒性能预测和模型选择提供依据。


\xsection{验证数据、计算模型与评价方法}{Validation Data, Computational Models, and Evaluation Methodology}

\xsubsection{公开基准数据与补充实验数据来源}{Sources of Public Benchmark Data and Supplementary Experimental Data}

本章验证数据主要包括公开基准数据和补充实验数据两部分。公开基准数据主要来源于 FRAPCON-4.0 Integral Assessment 报告中整理的燃料棒级实验数据库。该报告针对燃料中心温度、裂变气体释放、棒内自由体积、包壳腐蚀和环向应变等燃料性能评价量进行了系统评估，其中燃料中心温度和裂变气体释放是本文重点采用的两类公开验证数据。报告中的验证案例通常包含燃料棒几何参数、燃料类型、初始气隙、充气压力、功率历史、燃耗、边界条件以及实验测量结果，可为燃料性能程序的棒级验证提供基础数据。

考虑到本文第五章的研究重点为燃料中心温度和裂变气体释放行为，本章仅选取与上述两个评价量直接相关、且公开输入信息较完整的案例开展定量验证。对于报告附录中给出完整或基本完整 FRAPCON 输入卡的案例，本文将其中的燃料棒几何参数、轴向功率分布、功率历史、燃耗、初始充气条件和冷却剂边界条件整理并转换为本文程序所需的输入格式。对于 报告正文中虽给出实验说明或对比图，但未提供完整输入卡的案例，本文不将其列入定量计算样本，仅作为文献背景或结果讨论时的参考。

表~\ref{tab:frapcon_mox_steady_cases} 给出了 FRAPCON-4.0 综合评估数据库中与本文研究相关的稳态 MOX 燃料棒案例。表~\ref{tab:frapcon_mox_power_ramp_cases} 给出了 FRAPCON-4.0 综合评估数据库中 Beznau-1/Petten 短时间功率提升 MOX 燃料棒案例。该类燃料棒经历了基准辐照后的功率提升试验，实验数据包括棒平均燃耗、最高线功率、高功率持续时间和裂变气体释放结果，可用于评价裂变气体释放模型对短时间功率提升和高功率保持工况的响应能力。表中 ``$\checkmark$'' 表示该燃料棒在 FRAPCON-4.0 综合评估报告中被用于相应物理量的验证。大部分燃料棒间隙内填充的是纯氦气也有少部分填充氦—氙混合气体，用于评估间隙导热模型。寿期初（BOL)的线功率最高是60kW/m，寿期末（EOL）功率在快速上升的实验中最高可达47kW/m。燃料棒平均燃耗最高可达68GWd/MTU，但是对功率快速提高的燃料棒，最高的平均燃耗在47GWd/MTU。

\begin{xltabular}{\textwidth}{lXccc}
	\caption{FRAPCON-4.0 综合评估数据库中稳态 MOX 燃料棒案例}
	\label{tab:frapcon_mox_steady_cases}\\
	\toprule
	反应堆/试验装置 & 燃料棒 &
	棒平均燃耗/GWd$\cdot$MTU$^{-1}$ &
	温度-燃耗验证 & FGR 验证 \\
	\midrule
	\endfirsthead
	
	\multicolumn{5}{c}{续表~\thetable\quad FRAPCON-4.0 综合评估数据库中稳态 MOX 燃料棒案例}\\
	\toprule
	反应堆/试验装置 & 燃料棒 &
	棒平均燃耗/GWd$\cdot$MTU$^{-1}$ &
	温度-燃耗验证 & FGR 验证 \\
	\midrule
	\endhead
	
	\midrule
	\multicolumn{5}{r}{续下页}\\
	\endfoot
	
	\bottomrule
	\endlastfoot
	
	Halden HBWR & IFA-629-1r1 & 33 & $\checkmark$ &  \\
	Halden HBWR & IFA-629-1r2 & 29（FGR）/40（温度） & $\checkmark$ & $\checkmark$ \\
	
	Halden HBWR & IFA-610.2 & 56 & $\checkmark$ &  \\
	Halden HBWR & IFA-610.4 & 57 & $\checkmark$ &  \\
	
	Halden HBWR & IFA-648.1r1 & 62 & $\checkmark$ &  \\
	Halden HBWR & IFA-648.1r2 & 62 & $\checkmark$ &  \\
	
	Halden HBWR & IFA-629.3r5 & 72 & $\checkmark$ & $\checkmark$ \\
	Halden HBWR & IFA-629.3r6 & 68 & $\checkmark$ & $\checkmark$ \\
	
	Halden HBWR & IFA-606 Phase 2 & 49 & $\checkmark$ & $\checkmark$ \\
	
	Halden HBWR & IFA-651.1r1 & 22.41 & $\checkmark$ & $\checkmark$ \\
	Halden HBWR & IFA-651.1r3 & 21.73 & $\checkmark$ & $\checkmark$ \\
	Halden HBWR & IFA-651.1r6 & 20.27 & $\checkmark$ & $\checkmark$ \\
	
	ATR & PII C2 P5 & 21 &  & $\checkmark$ \\
	ATR & PIII C3 P6 & 30 &  & $\checkmark$ \\
	ATR & PIII C10 P13 & 30 &  & $\checkmark$ \\
	ATR & PIV C4 P7 & 40 &  & $\checkmark$ \\
	ATR & PIV C5 P8 & 50 &  & $\checkmark$ \\
	ATR & PIV C6 P9 & 50 &  & $\checkmark$ \\
	ATR & PIV C12 P15 & 50 &  & $\checkmark$ \\
	
	Gravelines-4 PWR & N06 & 48 &  & $\checkmark$ \\
	Gravelines-4 PWR & N12 & 57 &  & $\checkmark$ \\
	Gravelines-4 PWR & P16 & 53 &  & $\checkmark$ \\
	
	Halden HBWR & IFA-633.1r6 & 32 & $\checkmark$ & $\checkmark$ \\
	
	Beznau-1 & M504 H8 & 37.5 &  & $\checkmark$ \\
	Beznau-1 & M504 I2 & 43 &  & $\checkmark$ \\
	Beznau-1 & M504 K9 & 42.5 &  & $\checkmark$ \\
	Beznau-1 & M504 M9 & 44.2 &  & $\checkmark$ \\
	
	Beznau-1 & M308 Segment 2 & 57.5 &  & $\checkmark$ \\
	
	Halden HBWR & IFA-597.4/.5/.6/.7r10 & 35.7 & $\checkmark$ & $\checkmark$ \\
	Halden HBWR & IFA-597.4/.5/.6/.7r11 & 36.8 & $\checkmark$ & $\checkmark$ \\
	
	Fugen HBWR & E09 Rods Inner & 29.6 &  & $\checkmark$ \\
	Fugen HBWR & E09 Rods Intermediate & 39.3 &  & $\checkmark$ \\
	Fugen HBWR & E09 Rods Outer & 42 &  & $\checkmark$ \\
	
\end{xltabular}

\begin{table}[H]
	\centering
	\caption{FRAPCON-4.0 综合评估数据库中短时间功率提升 MOX 燃料棒案例}
	\label{tab:frapcon_mox_power_ramp_cases}
	\begin{tabularx}{\textwidth}{X c c c c}
		\toprule
		燃料棒 &
		棒平均燃耗/GWd$\cdot$MTU$^{-1}$ &
		最高线功率/kW$\cdot$m$^{-1}$ &
		高功率持续时间 & FGR 验证 \\
		\midrule
		
		\mbox{M501 HR-1} & 37 & 38.1 & 12 h & $\checkmark$ \\
		\mbox{M501 HR-2} & 37 & 35.7 & 12 h & $\checkmark$ \\
		\mbox{M501 HR-3} & 37 & 46.2 & 12 h & $\checkmark$ \\
		\mbox{M501 HR-4} & 36 & 47.0 & 12 h & $\checkmark$ \\
		
		\mbox{M501 MR-1} & 34 & 38.1 & 12 h & $\checkmark$ \\
		\mbox{M501 MR-2} & 34 & 41.9 & 12 h & $\checkmark$ \\
		\mbox{M501 MR-3} & 34 & 40.5 & 12 h & $\checkmark$ \\
		\mbox{M501 MR-4} & 33 & 41.7 & 20 min & $\checkmark$ \\
		
		\bottomrule
	\end{tabularx}
\end{table}

燃料中心温度验证主要选取 Halden 试验中的 MOX 燃料棒温度测量数据。这些燃料棒通常带有燃料中心温度测量装置，实验数据给出了中心温度随时间或燃耗的变化过程，可用于验证燃料热导率、间隙导热、燃耗退化以及裂变气体释放反馈对温度场预测的影响。根据公开输入卡的可获得性和本文实际建模情况，本章采用的中心温度公开验证案例包括 IFA-597-4/5/6/7r10、IFA-597-4/5/6/7r11、IFA-606、IFA-610.2、IFA-610.4、IFA-629.3r5、IFA-629.3r6、IFA-629-1r1、IFA-629-1r2、IFA-648r1、IFA-648r2、IFA-651-r1、IFA-651-r3 和 IFA-651-r6 等燃料棒。上述案例覆盖了不同燃耗水平、不同线功率历史以及不同实验组件条件下的 MOX 燃料中心温度演化过程。

需要说明的是，IFA-633.1r6 虽然在 FRAPCON-4.0 Integral Assessment 报告中被列为 MOX 中心温度和裂变气体释放评价案例，并且该燃料棒同时配置了燃料中心热电偶和压力传感器，但本文所使用的 FRAPCON-4.0 附录中未提供该案例可直接复现的完整输入卡。因此，为保证本文验证结果的可追溯性和输入参数的一致性，IFA-633.1r6 不纳入本章定量计算样本。该案例仅在必要时作为公开数据库覆盖范围的说明，不参与本文中心温度和 FGR 误差统计。

裂变气体释放验证案例进一步分为稳态辐照案例和短时间功率提升案例。稳态辐照案例主要用于验证模型在长时程稳态运行条件下对 EOL 裂变气体释放率的预测能力。本文采用的稳态 FGR 公开验证案例包括 ATR-PII-C2-P5、ATR-PIII-C3-P6、ATR-PIII-C10-P13、ATR-PIV-C4-P7、ATR-PIV-C5-P8、ATR-PIV-C6-P9、ATR-PIV-C12-P15、E09-inner、E09-intermediate、E09-outer、Gravelines-N6、Gravelines-N12、Gravelines-P16、IFA-597.4/5/6/7r10、IFA-597.4/5/6/7r11、IFA-606-P2、IFA-629.1、IFA-629.3r5、IFA-629.3r6、IFA-651-r1、IFA-651-r3、IFA-651-r6、M308-Segment2、M504-H8、M504-I2、M504-K9 和 M504-M9 等燃料棒。上述案例主要给出了棒平均燃耗、实验测得的 EOL FGR 以及相应的燃料棒运行历史，可用于评价 ForMas 和 SCIANTIX 模型在不同燃耗和功率水平下的裂变气体释放预测能力。

功率变化工况下的 FGR 验证主要采用 Beznau-1/Petten 相关短时间功率提升实验案例，包括 M501-HR-1、M501-HR-2、M501-HR-3、M501-HR-4、M501-MR-1、M501-MR-2、M501-MR-3 和 M501-MR-4。这些燃料棒经历了基准辐照后的功率提升试验，实验数据给出了棒平均燃耗、功率爬升终端线功率、保持时间以及测得的裂变气体释放率。该类案例可用于检验裂变气体模型对功率变化、温度升高和短时间高功率保持条件下气体释放增强现象的响应能力。

在数据使用过程中，本文基于 FRAPCON 附录中公开输入卡重新建立 MOOSE 框架下的燃料棒计算模型。对于燃料中心温度案例，主要提取实验测得的中心温度随时间或燃耗变化曲线数据，并与本文程序计算得到的燃料中心线平均温度进行对比。对于裂变气体释放案例，通过提取反应堆或棒内压反推得到的 FGR 数据，并与本文程序计算得到的总裂变气体释放分数进行对比。

除上述公开基准数据外，本文还有补充实验数据验证部分。补充实验数据主要用于考察模型在特定燃料体系、特定功率历史或公开数据库未覆盖工况下的适用性。对于非公开实验数据，本文仅给出开展计算所必需的燃料棒几何参数、初始气隙、充气压力、功率历史、燃耗、边界条件和实验测量量，并对燃料棒编号、实验装置来源及部分敏感设计参数进行脱敏处理。补充实验数据的验证方式与公开基准数据保持一致，即中心温度案例以燃料中心线温度为主要评价量，裂变气体释放案例以寿期末FGR、气体释放量或气体组分为主要评价量。

\xsubsection{ForMas 与 SCIANTIX 裂变气体模型设置}{Settings of the ForMas and SCIANTIX Fission Gas Models}

本章采用两类裂变气体行为模型开展燃料棒级验证，分别为 ForMas 模型和 SCIAN\allowbreak TIX 模型。两类模型的基本理论、控制方程、程序实现方式以及单元级验证结果已在第 4 章中进行了详细说明，因此本节不再重复模型公式推导，而重点说明二者在本章棒级验证中的计算设置、适用定位以及与温度场和间隙传热模型的耦合方式。

ForMas 模型来源于 Forsberg--Massih 裂变气体迁移理论，是一种面向燃料性能工程计算的简化机理模型。该模型将裂变气体释放过程主要描述为裂变气体在晶粒内生成、扩散、到达晶界、晶界累积以及达到饱和条件后释放的过程。根据第 4 章的模型实现，ForMas 模型在本文程序中以材料属性计算模块的形式嵌入 MOOSE 框架。其输入量主要包括局部温度、裂变率、燃耗、晶粒尺寸、时间步长以及模型相关经验参数；输出量主要包括当前时间步的裂变气体释放份额、累计释放份额以及产生份额。该模型采用解析解或半解析形式计算晶粒内扩散问题，避免了在有限元主方程中额外引入裂变气体浓度自由度，因此具有计算量小、数值稳定性好和便于长时程燃料棒性能计算的特点。

ForMas 模型中考虑的主要物理过程包括晶内扩散、辐照重溶、晶界气体累积以及饱和释放。当晶界处累积气体未达到释放阈值时，裂变气体主要存在晶内和晶界内；当晶界气体达到饱和条件后，部分气体被释放至燃料自由表面、芯块--包壳间隙或气腔中。释放后的晶界气体重新进入累积过程，因此该模型能够描述 ``累积--释放--再累积'' 的循环行为。对于稳态长期辐照案例，ForMas 模型可用于评价温度、燃耗和线功率水平对寿期末的 FGR 影响；对于短时间功率提升案例，该模型可用于考察温度升高引起的扩散增强和释放量增加现象。

SCIANTIX 模型是一种更高保真的细观机理裂变气体行为模型。公开文献中将 SCIANTIX 定义为面向燃料性能程序耦合的开源 0D 独立程序，其主要目标是在较低尺度机理计算与工程尺度燃料性能程序之间建立中间尺度模型。与 ForMas 模型相比，SCIANTIX 对裂变气体行为的描述更为细致，能够考虑晶内气体行为、晶界气体行为、气泡成核与长大、气泡重溶、气体释放、气体肿胀以及高燃耗结构形成等过程。新版 SCIANTIX 还进一步完善了程序结构、数值求解能力和分离效应验证数据库，可用于计算裂变气体释放和气体肿胀等工程量。

在本文程序中，SCIANTIX 通过动态链接库形式与 MOOSE 主程序进行耦合。MOOSE 主程序负责求解燃料棒尺度的温度场、燃耗场和裂变率场，并在每个时间步或指定耦合步向 SCIANTIX 传递局部温度、裂变率、燃耗、时间步长、初始晶粒尺寸以及相关模型选项。SCIANTIX 在细观尺度上完成裂变气体行为计算后，将裂变气体释放份额、晶内气体、晶界气体、气泡状态变量以及必要的气体肿胀信息返回至主程序。主程序随后将这些结果写回相应材料属性或后处理变量，用于更新燃料棒气腔中的 Xe、Kr 摩尔数、间隙气体组成以及后续时间步的热工边界条件。

两类裂变气体模型均与本文的间隙导热模型保持一致的耦合方式。具体而言，在当前时间步内，裂变气体模型首先根据温度、裂变率和燃耗等状态变量计算裂变气体释放量及其 Xe/Kr 组分；随后，GasMixture 后处理模块更新气隙和气腔中的 He、Xe、Kr 等气体物质的量与摩尔分数；在此基础上，GasGapConductance 根据混合气体组成、气体温度、气体压力和间隙宽度计算混合气体热导率和间隙换热系数；最终，GapHeatTransfer 利用更新后的间隙导热条件求解下一时间步的燃料温度场。新的温度场又会反馈至 ForMas 或 SCIANTIX 模型，进一步影响下一时间步的裂变气体扩散、累积与释放行为。通过该流程，本文实现了裂变气体释放、气体组分演化、间隙传热和燃料温度之间的双向耦合。

从计算定位上看，ForMas 模型主要用于工程尺度燃料棒长时程计算。其模型结构相对简单，参数数量较少，计算效率较高，适合用于大量公开基准案例、参数敏感性分析和燃料棒级性能趋势评价。SCIANTIX 模型则更适合用于需要细观机理解释的案例，尤其适用于分析晶内/晶界气体重新分配、高燃耗结构形成、气体肿胀以及复杂功率历史下的气体释放响应。由于 SCIANTIX 的输入变量和状态变量更多，对温度场、晶粒尺寸、时间步长和模型选项更为敏感，其计算成本和耦合复杂度高于 ForMas 模型。因此，本章在后续验证中分别给出 ForMas 和 SCIANTIX 的预测结果，并进一步比较二者在 FGR 预测精度、中心温度反馈、计算效率和适用范围方面的差异。

综上，ForMas 和 SCIANTIX 分别代表了工程化简化机理模型和细观高保真机理模型两种裂变气体释放建模路线。前者强调计算效率和燃料棒级长期稳定计算能力，后者强调对裂变气体细观演化过程的物理刻画能力。本章基于公开基准数据和补充实验数据对两类模型进行验证和对比，可为后续燃料性能分析中裂变气体模型的选择及温度--FGR--间隙导热耦合机制分析提供依据。

\xsubsection{计算输入、边界条件与结果提取方法}{Computational Inputs, Boundary Conditions, and Result Extraction Methods}

本节基于公开基准案例附录中给出的输入卡信息，建立本文程序所需的燃料棒计算模型。首先，根据文献输入卡中的燃料棒几何尺寸、燃料初始状态、功率历史、冷却剂边界条件和实验测量信息，将关键参数整理为统一格式的 text 文件。随后，利用 Python 脚本对这些基础数据进行预处理，自动生成 Cubit 建模脚本以及 MOOSE 输入卡中所需的几何参数和初始条件。预处理得到的主要输入量包括芯块外径、包壳内外径、初始气隙宽度、初始气腔压力、初始孔隙率、初始钚含量、初始 O/M 比、冷却剂入口温度、出口压力、质量流速、流道尺寸和水力直径等。

燃料棒几何模型采用二维轴对称 R--Z 形式建立，其中燃料芯块和包壳分别定义为 fuel 和 clad 计算区域。网格由 Cubit 生成，并导出为 Exodus 格式的 \texttt{.e} 文件。在 MOOSE 输入卡中，通过 \texttt{FileMeshGenerator} 读取该网格文件，并设置 \texttt{coord\_type = RZ} 进行轴对称计算。该处理方式能够在保留径向传热、轴向功率分布、芯块--包壳气隙和接触关系的同时，降低三维全尺寸燃料棒计算所需的网格规模和计算成本。

功率历史和功率分布均采用外部文件形式读入。棒平均线功率由
\texttt{rod\_\allowbreak ave\_\allowbreak lin\_\allowbreak pow.\allowbreak txt}
给出，并通过 \texttt{PiecewiseLinear} 函数
\texttt{rod\_\allowbreak ave\_\allowbreak lin\_\allowbreak pow\_\allowbreak fun}
读入；轴向功率分布和径向功率分布分别由
\texttt{axial\_\allowbreak power.\allowbreak txt}
和
\texttt{radial\_\allowbreak power.\allowbreak txt}
给出，并通过 \texttt{PiecewiseMultilinear} 函数读入。
在计算过程中，程序基于上述功率函数，通过
\texttt{Fission\allowbreak RateMOX}、\texttt{PowerDensityMOX}、
\texttt{EngineeringBurnupAux} 和 \texttt{BurnupAux}
计算局部裂变率、局部体积功率密度、工程燃耗和燃耗百分比。同时，\texttt{FastNeutronFlux} 材料模型根据功率水平计算快中子通量和中子注量，为包壳辐照生长和辐照蠕变模型提供输入。

热传导计算以燃料温度 \texttt{temperature(K)} 为主变量，控制方程由 \texttt{HeatConduction}、\texttt{HeatConductionTimeDerivative} 和 \texttt{NeutronHeatSource} 共同构成。其中，\texttt{NeutronHeat\allowbreak Source} 以裂变率为热源输入，实现裂变率到体热源的转换。燃料热物性由 \texttt{MAMOXThermal}、\texttt{MOXFuelDensity} 和 \texttt{MeltSolidTemperature} 等材料模型给出，考虑温度、孔隙率、钚浓度、O/M 比和燃耗对热导率、比热容、密度及熔点的影响。包壳热物性由 \texttt{Zr4Heat\allowbreak ConductionMaterial}、\texttt{Zr4Density} 和 \texttt{ZrPhaseTemperature} 等模型描述。

对于给出质量流量的条件，冷却剂边界采用包壳外表面对流换热形式。程序根据冷却剂入口温度、出口压力、质量流速和流道几何参数计算冷却剂状态。其中，\texttt{Water97\allowbreak FluidProperties} 提供水物性，\texttt{InletEnthalpyAux} 计算入口热力学状态$h$，\texttt{CoolantPerAux} 和 \texttt{LayeredSideIntegral} 用于计算沿轴向 $h$ 增量，\texttt{FluidPropertiesMaterialPHNEW} 根据压力和$h$计算局部冷却剂温度，\texttt{ComputeHw} 计算壁面换热系数。最终，包壳外表面的热边界条件通过 \texttt{Convective\allowbreak HeatFluxBC} 施加，其中环境温度为 \texttt{Tinf}，换热系数为 \texttt{Hw}。对于没有给出冷却剂质量的流量的工况，本文直接使用\texttt{DirichletBC}施加固定的包壳温度边界。冷却剂压力通过 \texttt{coolant\_pressure\_fun} 函数读入，并作为包壳外表面的力学压力边界条件。

裂变气体释放、气腔压力和气隙传热之间采用耦合计算。裂变气体释放由 \texttt{ForMas} 材料模型计算，其输入包括局部温度、裂变率、燃耗和气腔压力，输出为释放气体密度和气体生成量。释放后的 Xe/Kr 气体通过 \texttt{released\_gas\_moles} 后处理量积分得到，并传递给 \texttt{CavityPressure} 模块。气腔压力模型根据初始压力、初始气体摩尔数、释放气体摩尔数、气腔体积和气腔平均温度计算当前气腔压力 \texttt{plenum\_pressure}。其中，气腔体积由 \texttt{InternalVolume} 计算得到，并叠加弹簧区和缓冲区体积形成总自由体积 \texttt{volume\_total}；气腔平均温度由 \texttt{SideAverageValue} 在相关内表面上提取得到。

芯块--包壳间隙传热由 \texttt{GapHeatTransfer} 处理。该模型同时考虑气体导热、辐射换热和接触导热，并通过 \texttt{initial\_moles}、\texttt{initial\_gas\_types}、\texttt{plenum\_pressure} 和 \texttt{gas\_released} 接收气腔压力及释放气体信息。气体混合物状态由 \texttt{thermal\_\allowbreak contact\_\allowbreak gas\_\allowbreak mixture} 记录，并通过后处理量分别提取 He、Xe 和 Kr 的摩尔数。间隙换热系数以材料属性 \texttt{gap\_\allowbreak conductance} 的形式输出，并通过 \texttt{MaterialRealAux} 映射为辅助变量 \texttt{gap\_\allowbreak conductivity(W.m-2.K-1)}。因此，裂变气体释放会改变气隙气体组成和气腔压力，进而影响间隙热导率和燃料温度；新的温度场又会反馈至裂变气体释放模型，从而形成温度--裂变气体释放--间隙导热之间的耦合反馈。

力学计算采用准静态固体力学求解框架。燃料和包壳分别在 \texttt{SolidMechanics} 中定义，并以 \texttt{disp\_r} 和 \texttt{disp\_z} 为位移变量。燃料力学模型考虑热膨胀、重定位、致密化、固体裂变产物肿胀、气体肿胀和蠕变等本征应变或非弹性行为；包壳力学模型考虑热膨胀、辐照生长和蠕变行为。芯块--包壳之间的力学接触通过 \texttt{gap\_contact} 定义，采用无摩擦 penalty 接触形式。边界条件方面，燃料棒中心线施加径向位移约束，底部施加轴向位移约束，包壳外表面施加冷却剂压力边界。该力学模块用于获得芯块和包壳的应力、应变、位移和接触状态，并进一步影响间隙宽度、接触导热以及气腔体积变化。

图~\ref{fig:fuel_rod_multiphysics_flowchart} 给出了本文燃料棒多物理耦合计算流程。可以看出，程序以几何参数、功率历史和冷却剂边界为输入，通过功率--燃耗--中子通量模块驱动热传导、组分迁移、裂变气体释放和力学响应计算；温度场、裂变气体释放、气腔压力、间隙传热和力学接触之间通过材料属性、后处理量和边界条件实现数据传递，最终形成热--裂变气体--力学强耦合迭代过程。

\begin{figure}[!htb]
	\centering
	\includegraphics[width=\textwidth]{moose_fuel_rod_flowchart_editable.pdf}
	\caption{燃料棒级多物理耦合计算流程}
	\label{fig:fuel_rod_multiphysics_flowchart}
\end{figure}

\xsubsection{误差评价指标}{Error Evaluation Metrics}
\label{section:error_evalution_metrics}
本文主要针对燃料中心线温度和裂变气体释放率两个评价量进行误差分析。对于燃料中心线温度，实验数据通常随功率历史和燃耗变化而变化，因此本文在相同时间或相同燃耗位置提取程序计算得到的中心线平均温度，并与实验测得的中心线温度进行逐点对比。单个测点的温度绝对误差和相对误差分别定义为
\begin{equation}
	\Delta T_i = T_{\mathrm{cal},i}-T_{\mathrm{exp},i}
\end{equation}
\begin{equation}
	\delta T_i = \frac{T_{\mathrm{cal},i}-T_{\mathrm{exp},i}}{T_{\mathrm{exp},i}}\times 100\%
\end{equation}
式中：$T_{\mathrm{cal},i}$ 为第 $i$ 个测点对应的计算中心线平均温度，$T_{\mathrm{exp},i}$ 为实验测得的燃料中心线温度。

为评价不同燃料棒整体温度预测精度，进一步采用平均相对误差和均方根误差进行统计，其表达式分别为
\begin{equation}
	\overline{\delta T} =
	\frac{1}{N}\sum_{i=1}^{N}
	\left|
	\frac{T_{\mathrm{cal},i}-T_{\mathrm{exp},i}}{T_{\mathrm{exp},i}}
	\right|\times 100\%
\end{equation}
\begin{equation}
	\mathrm{RMSE}_T =
	\sqrt{
		\frac{1}{N}\sum_{i=1}^{N}
		\left(T_{\mathrm{cal},i}-T_{\mathrm{exp},i}\right)^2
	}
\end{equation}
式中：$N$ 为该燃料棒可用于对比的中心温度测点数量。

对于裂变气体释放行为，公开实验数据通常给出寿期末裂变气体释放率，因此本文采用寿期末计算得到的总裂变气体释放率与实验测量值进行对比。FGR 绝对误差定义为
\begin{equation}
	\Delta FGR = FGR_{\mathrm{cal}} - FGR_{\mathrm{exp}}
\end{equation}
式中：$FGR_{\mathrm{cal}}$ 为程序计算得到的寿期末裂变气体释放率，$FGR_{\mathrm{exp}}$ 为实验测得或由棒压反推得到的寿期末裂变气体释放率。

由于部分燃料棒的实验 FGR 数值较小，此时相对误差容易被放大，因此本文优先采用 FGR 绝对误差评价裂变气体释放模型的预测能力。在需要进行整体统计时，采用平均绝对误差表示为
\begin{equation}
	\overline{\Delta FGR} =
	\frac{1}{M}\sum_{j=1}^{M}
	\left|FGR_{\mathrm{cal},j}-FGR_{\mathrm{exp},j}\right|
\end{equation}
式中：$M$ 为参与 FGR 验证的燃料棒数量。对于实验 FGR 较高的案例，可进一步给出相对误差作为辅助评价指标：
\begin{equation}
	\delta FGR_j =
	\frac{FGR_{\mathrm{cal},j}-FGR_{\mathrm{exp},j}}
	{FGR_{\mathrm{exp},j}}\times 100\%
\end{equation}

综上，中心线温度验证主要关注随时间或燃耗变化过程中的逐点温度偏差及整体统计误差；裂变气体释放验证主要关注寿期末 FGR 的计算值与实验值差异。上述评价指标为后续不同燃料棒案例以及 ForMas 和 SCIANTIX 两类裂变气体模型的预测能力对比提供统一依据。

\xsection{燃料中心温度验证}{Validation of Fuel Centerline Temperature}

\xsubsection{随燃耗演化的中心温度验证}{Validation of Centerline Temperature Evolution with Burnup}

为验证所建立燃料性能分析程序对燃料棒热行为的预测能力，本文选取 IFA 系列 MOX 燃料棒的中心温度实验数据开展对比分析。输入卡中给出了相应燃料棒的功率历史、冷却剂边界条件及轴向功率分布等运行参数，但公开实验数据的完整程度存在差异：部分燃料棒仅给出了若干离散燃耗点处的中心温度结果，而少数燃料棒提供了覆盖较长燃耗区间甚至全寿期的温度演化曲线。因此，本文首先利用所有可获得的实验数据对模型进行总体趋势验证，随后重点选取具有连续燃耗演化数据的燃料棒开展误差来源分析。

图~ 给出了不同 MOX 燃料棒中心温度随燃耗演化的模拟结果与实验结果对比。总体来看，程序能够较好地再现实验曲线的主要变化趋势，包括功率变化引起的中心温度升降、燃耗积累导致的热物性退化影响，以及间隙闭合或燃料-包壳接触后中心温度降低的现象。不同燃料棒之间的温度水平差异主要来源于线功率历史、轴向功率分布、燃料几何尺寸、初始间隙以及冷却剂温度边界条件的差异。


\begin{figure}[!htb]
	\centering
	
	\begin{subfigure}[t]{0.48\textwidth}
		\centering
		\includegraphics[width=\linewidth]{IFA-597-4-5-6-7r10.pdf}
		\caption{IFA-597-4/.5/.6/.7r10}
		\label{fig:temp_ifa597_r10}
	\end{subfigure}
	\hfill
	\begin{subfigure}[t]{0.48\textwidth}
		\centering
		\includegraphics[width=\linewidth]{IFA-597-4-5-6-7r11.pdf}
		\caption{IFA-597-4/.5/.6/.7r11}
		\label{fig:temp_ifa597_r11}
	\end{subfigure}
	
	\vspace{0.5em}
	
	\begin{subfigure}[t]{0.48\textwidth}
		\centering
		\includegraphics[width=\linewidth]{IFA-606.pdf}
		\caption{IFA-606 Phase 2}
		\label{fig:temp_ifa606}
	\end{subfigure}
	
	\caption{IFA-597 和 IFA-606 燃料棒中心温度随燃耗演化验证结果}
	\label{fig:temp_validation_ifa597_ifa606}
\end{figure}

图~\ref{fig:temp_validation_ifa597_ifa606} 给出了 IFA-597-4/.5/.6/.7r10、IFA-597-4/.5/.6/.7r11 以及 IFA-606 Phase 2 的中心温度验证结果。可以看出，模拟曲线能够较好再现实验中心温度随燃耗变化的总体趋势。对于 IFA-597 系列燃料棒，中心温度随功率历史变化表现出明显波动，模拟结果能够捕捉主要升降趋势，说明程序对线功率变化和燃料热响应的一阶行为具有较好的描述能力。对于 IFA-606 Phase 2，温度曲线在中后期出现较明显的下降趋势，这与燃料-包壳间隙逐渐闭合、间隙导热增强以及接触传热作用增强有关。

\begin{figure}[!htb]
	\centering
	
	\begin{subfigure}[t]{0.48\textwidth}
		\centering
		\includegraphics[width=\linewidth]{IFA-610-2.pdf}
		\caption{IFA-610.2}
	\end{subfigure}
	\hfill
	\begin{subfigure}[t]{0.48\textwidth}
		\centering
		\includegraphics[width=\linewidth]{IFA-610-4.pdf}
		\caption{IFA-610.4}
	\end{subfigure}
	
	\vspace{0.5em}
	
	\begin{subfigure}[t]{0.48\textwidth}
		\centering
		\includegraphics[width=\linewidth]{IFA-629-1r1.pdf}
		\caption{IFA-629-1r1}
	\end{subfigure}
	\hfill
	\begin{subfigure}[t]{0.48\textwidth}
		\centering
		\includegraphics[width=\linewidth]{IFA-629-1r2.pdf}
		\caption{IFA-629-1r2}
	\end{subfigure}
	
	\caption{IFA-610 与 IFA-629-1 燃料棒中心温度验证结果}
	\label{fig:temp_validation_ifa610_ifa6291}
\end{figure}

图~\ref{fig:temp_validation_ifa610_ifa6291} 对比了 IFA-610 与 IFA-629-1 系列燃料棒的中心温度。该组燃料棒覆盖约 29--57 GWd/MTM 的燃耗范围，能够反映中等燃耗区间内燃料热导率退化与间隙演化的耦合作用。模拟结果总体上能够跟随实验温度的变化趋势，但在部分燃耗区间仍存在一定偏差。该偏差可能与轴向功率分布、冷却剂温度边界条件、燃料重定位模型以及燃料-包壳接触发生时间有关。

\begin{figure}[!htb]
	\centering
	
	\begin{subfigure}[t]{0.48\textwidth}
		\centering
		\includegraphics[width=\linewidth]{IFA-648-r1.pdf}
		\caption{IFA-648-r1}
	\end{subfigure}
	\hfill
	\begin{subfigure}[t]{0.48\textwidth}
		\centering
		\includegraphics[width=\linewidth]{IFA-648-r2.pdf}
		\caption{IFA-648-r2}
	\end{subfigure}
	
	\vspace{0.5em}
	
	\begin{subfigure}[t]{0.48\textwidth}
		\centering
		\includegraphics[width=\linewidth]{IFA-629.3r5.pdf}
		\caption{IFA-629.3r5}
	\end{subfigure}
	\hfill
	\begin{subfigure}[t]{0.48\textwidth}
		\centering
		\includegraphics[width=\linewidth]{IFA-629.3r6.pdf}
		\caption{IFA-629.3r6}
	\end{subfigure}
	
	\caption{IFA-648-r 与 IFA-629.3r 燃料棒中心温度验证结果}
	\label{fig:temp_validation_high_burnup}
\end{figure}

图~\ref{fig:temp_validation_high_burnup} 给出了 IFA-648 与 IFA-629.3 高燃耗燃料棒的中心温度验证结果。相比中低燃耗燃料棒，高燃耗阶段中心温度受燃料热导率退化、裂变气体释放、燃料肿胀以及燃料-包壳接触状态的共同影响更为明显。模拟曲线能够再现实验温度的主要变化趋势，说明模型能够在较宽燃耗范围内描述 MOX 燃料的热行为。然而，高燃耗区间内局部偏差有所增大，表明燃料热物性退化和间隙传热模型仍是影响中心温度预测精度的重要因素。
\begin{figure}[!htb]
	\centering
	
	\begin{subfigure}[t]{0.48\textwidth}
		\centering
		\includegraphics[width=\linewidth]{IFA-651-r1.pdf}
		\caption{IFA-651-r1}
		\label{fig:temp_ifa_651_r1}
	\end{subfigure}
	\hfill
	\begin{subfigure}[t]{0.48\textwidth}
		\centering
		\includegraphics[width=\linewidth]{IFA-651-r3.pdf}
		\caption{IFA-651-r3}
		\label{fig:temp_ifa_651_r3}
	\end{subfigure}
	
	\vspace{0.5em}
	
	\begin{subfigure}[t]{0.48\textwidth}
		\centering
		\includegraphics[width=\linewidth]{IFA-651-r6.pdf}
		\caption{IFA-651-r6}
		\label{fig:temp_ifa_651_r6}
	\end{subfigure}
	
	\caption{IFA-651 燃料棒中心温度随燃耗演化验证结果}
	\label{fig:temp_validation_ifa651r6}
\end{figure}

图~\ref{fig:temp_validation_ifa651r6} 展示了 IFA-651 系列三根燃料棒的中心温度验证结果。该系列燃料棒燃耗相对较低，但功率历史和边界条件变化较明显，因此可用于考察程序对运行工况变化的响应能力。模拟结果能够较好反映中心温度随燃耗和功率变化的演化趋势，说明所建立模型对 MOX 燃料棒稳态热响应具有较好的适用性。

在多数燃耗区间内，模拟曲线能够较好跟随实验曲线的变化趋势，表明程序对燃料中心温度的一阶响应具有较好的预测能力。值得注意的是，模拟结果显示，在当前轴向约束条件下，燃料芯块的轴向膨胀受到底部固定边界的限制，热膨胀、辐照肿胀、重定位以及密实化等本征应变共同改变燃料与包壳之间的间隙状态。随着燃耗增加和功率历史变化，部分轴向位置首先发生燃料-包壳接触。接触发生后，燃料表面与包壳内表面之间的等效间隙热阻显著降低，间隙导热能力增强，因此在相同线功率条件下燃料中心温度出现下降。该现象说明中心温度不仅受燃料热导率和线功率控制，也受到燃料-包壳机械接触状态的显著影响。在本算例中，接触首先出现在燃料顶部附近，这与轴向约束、温度场分布以及局部变形累积有关。由于顶部区域的轴向自由变形和径向膨胀耦合更为明显，局部间隙更容易闭合，从而使该区域的间隙导热提前增强。

为定量评价不同燃料棒中心温度预测精度，本文在实验曲线对应的燃耗或时间位置对计算结果进行插值，得到与实验测点一一对应的计算温度。对于具有时间信息的实验曲线，首先采用与图中一致的峰值对齐方法对实验时间进行平移，然后在对齐后的实验时间位置提取计算温度；对于仅给出燃耗坐标的实验曲线，则在统一燃耗单位后按燃耗位置进行插值。误差采用\ref{section:error_evalution_metrics}节给出的计算方法，统计结果见表~\ref{tab:centerline_temperature_error}。

\begin{table}[htbp]
	\centering
	\caption{MOX燃料棒中心温度验证误差统计}
	\label{tab:centerline_temperature_error}
	\begin{tabular}{lccccc}
		\toprule
		燃料棒 & 燃耗范围/GWd$\cdot$MTM$^{-1}$ & 对比点数 & MBE/K & $\overline{\delta T}$/\% & RMSE/K \\
		\midrule
		IFA-597-4/.5/.6/.7r10 & 0.55--36.10 & 520  & -71.8  & 8.71  & 131.7 \\
		IFA-597-4/.5/.6/.7r11 & 0.41--34.86 & 1419 & 10.7   & 13.37 & 175.4 \\
		IFA-606                & 49.33--49.83 & 239  & -132.2 & 12.94 & 205.4 \\
		IFA-610.2              & 57.24--60.52 & 598  & -32.6  & 4.82  & 52.0 \\
		IFA-610.4              & 59.55--61.22 & 606  & 9.3    & 1.59  & 19.7 \\
		IFA-629.3r5            & 73.17--77.67 & 467  & 13.2   & 7.54  & 108.4 \\
		IFA-629.3r6            & 69.08--73.52 & 461  & -18.7  & 9.94  & 112.8 \\
		IFA-629-1r1            & 30.52--38.61 & 358  & 140.6  & 24.43 & 291.4 \\
		IFA-629-1r2            & 30.39--37.55 & 1920 & 236.5  & 40.30 & 355.9 \\
		IFA-648r1              & 53.05--63.28 & 493  & 40.7   & 5.61  & 60.2 \\
		IFA-648r2              & 52.69--63.99 & 441  & 118.3  & 13.83 & 130.1 \\
		IFA-651-r1             & 0.40--24.50  & 800  & -87.0  & 14.73 & 227.8 \\
		IFA-651-r3             & 0.10--23.13  & 457  & -60.0  & 6.70  & 147.4 \\
		IFA-651-r6             & 0.09--21.99  & 495  & -249.9 & 21.98 & 299.6 \\
		\bottomrule
	\end{tabular}
\end{table}

从表中可以看出，不同燃料棒中心温度预测误差存在一定差异。整体而言，多数燃料棒的平均相对误差处于较低水平，其中 IFA-610.2、IFA-610.4、IFA-648r1、IFA-651-r3 等燃料棒的平均相对误差小于 7\%，表明程序能够较好地预测 MOX 燃料棒中心温度随燃耗和功率历史变化的总体趋势。IFA-610.4 的 RMSE 仅为 19.7 K，平均相对误差为 1.59\%，说明该算例中燃料热导率、间隙传热和冷却剂边界条件的综合作用与实验结果较为一致。

从平均偏差 MBE 的符号来看，不同燃料棒存在一定的系统性偏差。IFA-629-1r1、IFA-629-1r2 和 IFA-648r2 的 MBE 为正，说明这些算例中程序整体倾向于高估中心温度；而 IFA-606、IFA-651-r1、IFA-651-r3 和 IFA-651-r6 的 MBE 为负，说明这些算例中程序整体偏低估中心温度。该差异表明中心温度偏差并非由单一因素控制，而是由燃料热导率演化、燃料-包壳间隙演化、接触传热、冷却剂温度边界以及功率历史共同决定。


\xsubsection{中心温度预测偏差来源分析}{Analysis of Deviation Sources in Centerline Temperature Prediction}
\subsubsection{缩放因子敏感性分析方法}
为进一步识别中心温度预测偏差的来源，本文选取具有较完整寿期温度数据的五根燃料棒（IFA-597-4/.5/.6/.7r10 IFA-597-4/.5/.6/.7r11 IFA-651-r1 IFA-651-r3 IFA-651-r6）开展敏感性分析。由于燃料中心温度由体积释热、燃料热导率、燃料-包壳间隙传热、冷却剂边界条件以及燃料-包壳接触状态共同决定，仅通过单一温度曲线难以直接判断误差来源。因此，本文引入缩放因子方法，对主要热工和力学控制参数进行扰动分析，并比较不同扰动条件下模拟曲线相对于实验曲线的变化方向和变化幅度。

设基准模型计算结果为 \(T_\mathrm{base}\)，参数缩放后的计算结果为 \(T_\alpha\)，则温度响应可表示为
\[
\Delta T_\alpha(B)
=
T_\alpha(B)-T_\mathrm{base}(B),
\]
其中 \(B\) 为燃耗。主要依据两个准则。第一，若某一参数扰动后计算曲线整体向实验曲线靠近，且 RMSE 或平均相对误差明显降低，则说明该参数相关模型可能是误差的重要来源。第二，若某一参数扰动引起的温度响应主要出现在特定燃耗阶段，则说明该误差来源具有阶段性。


根据中心温度的主要控制机制，本文重点分析以下三类缩放因子的影响：

\begin{itemize}
	\item 燃料热导率缩放因子 \(\alpha_k\)。
	燃料热导率直接控制燃料内部径向温降。在相同线功率条件下，降低燃料热导率会增大燃料中心至燃料表面的温度梯度，使中心温度升高；提高燃料热导率则会降低中心温度。因此，当基准模型在高燃耗区间系统性低估中心温度时，可能说明燃料热导率退化不足；反之，若基准模型系统性高估中心温度，则可能说明燃料热导率偏低，或间隙传热能力被低估。
	
	\item 燃料-包壳间隙导热缩放因子 \(\alpha_h\)。
	燃料-包壳间隙等效导热系数决定燃料表面至包壳内表面的热阻，是影响中心温度的重要因素。提高间隙导热系数会降低燃料表面温度，并进一步降低燃料中心温度；降低间隙导热系数则会使中心温度升高。该因素在寿期早期通常更为敏感，因为此时燃料-包壳之间仍存在明显气隙，热阻主要由气体导热、辐射换热以及表面粗糙度共同决定。随着燃耗增加和燃料-包壳接触发生，间隙热阻降低，中心温度对 \(h_\mathrm{gap}\) 的敏感性逐渐减弱。
	
	\item 燃料几何演化缩放因子 \(\alpha_\varepsilon\)。
	燃料重定位、密实化和辐照肿胀通过改变燃料几何尺寸和燃料-包壳间隙状态间接影响中心温度。重定位和肿胀通常促进间隙闭合，使间隙热阻降低，从而降低中心温度；密实化则可能增大燃料-包壳间隙，使间隙热阻增加并导致中心温度升高。因此，这类模型对中心温度的影响不一定表现为简单的单调变化，而是主要体现为间隙闭合时间、接触压力大小以及接触区域分布的改变。
\end{itemize}

为定量分析不同物理模型对中心温度的影响，本文分别引入燃料热导率缩放因子 \(\alpha_k\)、间隙导热缩放因子 \(\alpha_h\) 和燃料变形缩放因子 \(\alpha_\varepsilon\)，其作用形式可表示为
\begin{align}
	k_f^\ast &= \alpha_k k_f \nonumber \\
	h_\mathrm{gap}^\ast &= \alpha_h h_\mathrm{gap} \nonumber \\
	\varepsilon^\ast_\mathrm{fuel}
	&=
	\alpha_\varepsilon
	\left(
	\varepsilon_\mathrm{reloc}
	+
	\varepsilon_\mathrm{dens}
	+
	\varepsilon_\mathrm{swelling}
	\right) \label{eq:fuel_params}
\end{align}
式中：\(k_f\) 为燃料热导率，\(h_\mathrm{gap}\) 为燃料-包壳间隙等效导热系数，\(\varepsilon_\mathrm{fuel}\) 表示与燃料几何演化相关的等效本征应变。

采用单因素扰动方法，在保持其他模型参数不变的条件下，缩放因子分别取
\begin{align}
	\alpha_k &= 0.9,\ 1.0,\ 1.1 \nonumber \\
	\alpha_h &= 0.5,\ 1.0,\ 2.0 \nonumber \\
	\alpha_\varepsilon &= 0.5,\ 1.0,\ 1.5 \label{eq:alpha_params}
\end{align}

\subsubsection{代表燃料棒敏感性分析}

通过比较不同扰动条件下中心温度曲线相对于实验曲线的变化方向和误差指标变化，可以进一步识别中心温度预测偏差对不同物理模型的敏感性。为说明缩放因子分析的具体作用，本文首先以 IFA-597-4-5-6-7r11 为代表燃料棒开展详细分析。该燃料棒的基准计算结果整体高于实验值，因此适合用于分析增强间隙传热、提高燃料热导率以及增强燃料几何演化对中心温度高估偏差的修正作用。

\begin{figure}[!htb]
	\centering
	
	\begin{subfigure}[t]{0.48\textwidth}
		\centering
		\includegraphics[width=\linewidth]{IFA-597-4-5-6-7r11/IFA-597-4-5-6-7r11_temperature_fuel_thermal_conductivity_alpha_comparison.pdf}
		\caption{燃料热导率缩放因子}
		\label{fig:ifa597r11_temp_alpha_fuel_thermal_conductivity}
	\end{subfigure}
	\hfill
	\begin{subfigure}[t]{0.48\textwidth}
		\centering
		\includegraphics[width=\linewidth]{IFA-597-4-5-6-7r11/IFA-597-4-5-6-7r11_temperature_gap_conductivity_alpha_comparison.pdf}
		\caption{间隙传热缩放因子}
		\label{fig:ifa597r11_temp_alpha_gap_conductance}
	\end{subfigure}
	
	\vspace{0.5em}
	
	\begin{subfigure}[t]{0.48\textwidth}
		\centering
		\includegraphics[width=\linewidth]{IFA-597-4-5-6-7r11/IFA-597-4-5-6-7r11_temperature_geometricstructure_alpha_comparison.pdf}
		\caption{燃料几何演化缩放因子}
		\label{fig:ifa597r11_temp_alpha_geometric_evolution}
	\end{subfigure}
	
	\caption{IFA-597-4-5-6-7r11 燃料中心温度对不同物理模型缩放因子的响应}
	\label{fig:ifa597r11_temperature_alpha_sensitivity}
\end{figure}

图~\ref{fig:ifa597r11_temperature_alpha_sensitivity} 给出了 IFA-597-4-5-6-7r11 燃料棒中心温度对不同物理模型缩放因子的响应。为进一步定量比较不同缩放因子对中心温度预测误差的影响，表~\ref{tab:ifa597r11_temperature_alpha_sensitivity_error} 给出了对应的误差统计结果。其中，\(\alpha=1.0\) 表示基准模型，因此在三类扰动分析中均作为对照工况列出。

\begin{table}[htbp]
	\centering
	\caption{IFA-597-4-5-6-7r11 不同缩放因子下中心温度预测误差统计}
	\label{tab:ifa597r11_temperature_alpha_sensitivity_error}
	\small
	\begin{tabular*}{\textwidth}{
			@{\extracolsep{\fill}}
			llccccc
			@{}
		}
		\toprule
		扰动对象 & 缩放因子 & 燃耗范围 & 对比点数 & MBE/K & $\overline{\delta T}$/\% & RMSE/K \\
		\midrule
		力学 & $\alpha_{\varepsilon}=0.5$ & 0.41--34.73 & 1414 & 341.0 & 35.57 & 390.4 \\
		力学 & $\alpha_{\varepsilon}=1.0$ & 0.41--34.86 & 1419 & 177.8 & 20.04 & 258.0 \\
		力学 & $\alpha_{\varepsilon}=1.5$ & 0.41--35.07 & 1426 & 54.0 & 13.80 & 185.6 \\
		\midrule
		间隙传热 & $\alpha_h=0.5$ & 0.41--35.00 & 1424 & 190.6 & 21.11 & 271.7 \\
		间隙传热 & $\alpha_h=1.0$ & 0.41--34.86 & 1419 & 177.8 & 20.04 & 258.0 \\
		间隙传热 & $\alpha_h=2.0$ & 0.41--34.73 & 1414 & 101.9 & 15.35 & 211.6 \\
		\midrule
		燃料热导率 & $\alpha_k=0.9$ & 0.41--34.93 & 1422 & 189.8 & 20.88 & 270.1 \\
		燃料热导率 & $\alpha_k=1.0$ & 0.41--34.86 & 1419 & 177.8 & 20.04 & 258.0 \\
		燃料热导率 & $\alpha_k=1.1$ & 0.41--34.79 & 1417 & 136.8 & 17.08 & 225.7 \\
		\bottomrule
	\end{tabular*}
\end{table}

由表~\ref{tab:ifa597r11_temperature_alpha_sensitivity_error} 可以看出，在基准模型下，IFA-597-4-5-6-7r11 的中心温度预测结果整体高于实验值，MBE 为 177.8 K，平均相对误差为 20.04\%，RMSE 为 258.0 K。不同缩放因子对误差指标的影响程度存在明显差异，说明该燃料棒中心温度偏差并非由单一物理过程控制，而是由燃料热导率、间隙传热和燃料几何演化共同作用的结果。

首先，燃料几何演化缩放因子对中心温度预测结果影响最为显著。当 \(\alpha_{\varepsilon}\) 由 0.5 增大至 1.5 时，MBE 由 341.0 K 降低至 54.0 K，平均相对误差由 35.57\% 降低至 13.80\%，RMSE 由 390.4 K 降低至 185.6 K，最大绝对偏差也由 910.6 K 降低至 705.1 K。这说明增强燃料重定位、密实化和肿胀等几何演化效应后，计算得到的中心温度整体向实验曲线靠近。其物理原因在于，燃料几何演化会改变燃料-包壳间隙状态和接触发生时间。较大的几何演化缩放因子会促进间隙闭合和接触传热增强，从而降低燃料中心温度。因此，基准模型对中心温度的高估可能与燃料几何演化效应偏弱、间隙闭合不足或接触传热建立偏晚有关。

其次，间隙传热缩放因子对中心温度也具有明显影响。当 \(\alpha_h\) 从 0.5 增大至 2.0 时，MBE 由 190.6 K 降低至 101.9 K，平均相对误差由 21.11\% 降低至 15.35\%，RMSE 由 271.7 K 降低至 211.6 K。该结果表明，增强燃料-包壳间隙传热能力可以有效降低中心温度，并改善与实验结果的符合程度。因此，基准模型中部分温度高估也可能来源于间隙热阻偏大，尤其是在尚未完全接触的燃耗区间，间隙气体导热、辐射换热和表面粗糙度等因素会显著影响燃料中心温度。

相比之下，燃料热导率缩放因子的影响相对较弱但仍然明显。当 \(\alpha_k\) 从 0.9 增大至 1.1 时，MBE 由 189.8 K 降低至 136.8 K，平均相对误差由 20.88\% 降低至 17.08\%，RMSE 由 270.1 K 降低至 225.7 K。提高燃料热导率会降低燃料内部径向温降，从而降低中心温度。该趋势说明燃料热导率模型对中心温度偏差具有一定贡献，但其改善幅度小于燃料几何演化和间隙传热缩放因子。

综合来看，对于 IFA-597-4-5-6-7r11 燃料棒，中心温度预测偏差主要受燃料几何演化和间隙传热模型控制，燃料热导率模型也具有一定影响。特别是 \(\alpha_{\varepsilon}=1.5\) 工况下误差指标改善最为明显，表明燃料-包壳间隙状态和接触传热演化是该算例中心温度偏差的重要来源。需要说明的是，缩放因子分析并非用于重新标定模型参数，而是用于识别不同物理过程对中心温度预测偏差的敏感性。后续模型改进应重点关注燃料重定位、密实化、肿胀以及间隙传热模型对接触建立过程和热阻演化的影响。

\subsubsection{不同燃料棒敏感性差异}
需要指出的是，IFA-597-4-5-6-7r11 仅代表基准模型高估中心温度的一类算例。对于其他燃料棒，基准模型相对于实验结果的偏差方向和燃耗阶段分布可能不同，因此相同缩放因子对误差指标的影响也可能发生变化。为避免单一算例结论的偶然性，本文进一步对 IFA-597-4-5-6-7r10、IFA-651-r1、IFA-651-r3 和 IFA-651-r6 开展相同的缩放因子分析，并比较不同燃料棒之间的敏感性差异。

\begin{figure}[!htb]
	\centering
	
	\begin{subfigure}[t]{0.48\textwidth}
		\centering
		\includegraphics[width=\linewidth]{IFA-597-4-5-6-7r10/IFA-597-4-5-6-7r10_temperature_fuel_thermal_conductivity_alpha_comparison.pdf}
		\caption{燃料热导率缩放因子}
		\label{fig:ifa597r10_temp_alpha_fuel_thermal_conductivity}
	\end{subfigure}
	\hfill
	\begin{subfigure}[t]{0.48\textwidth}
		\centering
		\includegraphics[width=\linewidth]{IFA-597-4-5-6-7r10/IFA-597-4-5-6-7r10_temperature_gap_conductivity_alpha_comparison.pdf}
		\caption{间隙传热缩放因子}
		\label{fig:ifa597r10_temp_alpha_gap_conductance}
	\end{subfigure}
	
	\vspace{0.5em}
	
	\begin{subfigure}[t]{0.48\textwidth}
		\centering
		\includegraphics[width=\linewidth]{IFA-597-4-5-6-7r10/IFA-597-4-5-6-7r10_temperature_geometricstructure_alpha_comparison.pdf}
		\caption{燃料几何演化缩放因子}
		\label{fig:ifa597r10_temp_alpha_geometric_evolution}
	\end{subfigure}
	
	\caption{IFA-597-4-5-6-7r10 燃料中心温度对不同物理模型缩放因子的响应}
	\label{fig:ifa597r10_temperature_alpha_sensitivity}
\end{figure}

\begin{table}[H]
	\centering
	\caption{IFA-597-4-5-6-7r10不同缩放因子下中心温度预测误差统计}
	\label{tab:ifa597r10_temperature_alpha_sensitivity_error}
	\small
	\begin{tabular*}{0.90\textwidth}{
			@{\extracolsep{\fill}}
			llccccc
			@{}
		}
		\toprule
		扰动对象 & 缩放因子 & 燃耗范围 & 对比点数 & MBE/K & $\overline{\delta T}$/\% & RMSE/K \\
		\midrule
		力学 & $\alpha_{\varepsilon}=0.5$ & 0.55--36.10 & 520 & 81.1 & 8.71 & 136.1 \\
		力学 & $\alpha_{\varepsilon}=1.0$ & 0.55--36.10 & 520 & -71.8 & 8.71 & 131.7 \\
		力学 & $\alpha_{\varepsilon}=1.5$ & 0.55--36.10 & 520 & -167.8 & 14.58 & 201.7 \\
		\midrule
		间隙传热 & $\alpha_h=0.5$ & 0.55--36.10 & 520 & -72.0 & 9.92 & 151.3 \\
		间隙传热 & $\alpha_h=1.0$ & 0.55--36.10 & 520 & -71.8 & 8.71 & 131.7 \\
		间隙传热 & $\alpha_h=2.0$ & 0.55--36.10 & 520 & -117.8 & 10.90 & 152.1 \\
		\midrule
		燃料热导率 & $\alpha_k=0.9$ & 0.55--36.10 & 520 & -60.9 & 9.38 & 140.2 \\
		燃料热导率 & $\alpha_k=1.0$ & 0.55--36.10 & 520 & -71.8 & 8.71 & 131.7 \\
		燃料热导率 & $\alpha_k=1.1$ & 0.55--36.10 & 520 & -100.5 & 9.90 & 143.7 \\
		\bottomrule
	\end{tabular*}
\end{table}

图~\ref{fig:ifa597r10_temperature_alpha_sensitivity} 和表~\ref{tab:ifa597r10_temperature_alpha_sensitivity_error} 给出了 IFA-597-4-5-6-7r10 的敏感性分析结果。与 r11 不同，该燃料棒基准模型整体低估中心温度，MBE 为 \(-71.8\) K。对于燃料几何演化缩放因子，降低 \(\alpha_{\varepsilon}\) 可使 MBE 由负值转为正值，但 RMSE 并未明显降低；增大 \(\alpha_{\varepsilon}\) 则进一步降低中心温度并增大误差。间隙传热和燃料热导率缩放因子的影响相对较弱，基准模型在 RMSE 和平均相对误差上已经接近最优。该结果说明，r10 的温度偏差方向与 r11 不同，缩放因子分析应结合基准模型的高估或低估特征进行判断。

% =========================
% 第一页：子图 (a)、(b)
% =========================
\begin{figure}[!htb]
	\centering
	
	\begin{subfigure}[t]{0.48\textwidth}
		\centering
		\includegraphics[width=\linewidth]
		{IFA-651-r1/IFA-651-r1_temperature_fuel_thermal_conductivity_alpha_comparison.pdf}
		\caption{燃料热导率缩放因子}
		\label{fig:ifa651r1_temp_alpha_fuel_thermal_conductivity}
	\end{subfigure}
	\hfill
	\begin{subfigure}[t]{0.48\textwidth}
		\centering
		\includegraphics[width=\linewidth]
		{IFA-651-r1/IFA-651-r1_temperature_gap_conductivity_alpha_comparison.pdf}
		\caption{间隙传热缩放因子}
		\label{fig:ifa651r1_temp_alpha_gap_conductance}
	\end{subfigure}
	
	\caption{IFA-651-r1 燃料中心温度对不同物理模型缩放因子的响应}
	\label{fig:ifa651r1_temperature_alpha_sensitivity}
\end{figure}


% 如果你明确要求下面这个图从下一页开始，可以加
\clearpage


% =========================
% 第二页：子图 (c)
% =========================
\begin{figure}[!htb]
	\ContinuedFloat
	\centering
	
	\begin{subfigure}[t]{0.48\textwidth}
		\centering
		\includegraphics[width=\linewidth]
		{IFA-651-r1/IFA-651-r1_temperature_geometricstructure_alpha_comparison.pdf}
		\caption{燃料几何演化缩放因子}
		\label{fig:ifa651r1_temp_alpha_geometric_evolution}
	\end{subfigure}
	
	\caption{IFA-651-r1 燃料中心温度对不同物理模型缩放因子的响应（续）}
\end{figure}

\begin{table}[H]
	\centering
	\caption{IFA-651-r1 不同缩放因子下中心温度预测误差统计}
	\label{tab:ifa651r1_temperature_alpha_sensitivity_error}
	\small
	\begin{tabular*}{0.9\textwidth}{
			@{\extracolsep{\fill}}
			llccccc
			@{}
		}
		\toprule
		扰动对象 & 缩放因子 & 燃耗范围 & 对比点数 & MBE/K & $\overline{\delta T}$/\% & RMSE/K \\
		\midrule
		力学 & $\alpha_{\varepsilon}=0.5$ & 0.40--24.45 & 799 & -18.3 & 12.20 & 192.3 \\
		力学 & $\alpha_{\varepsilon}=1.0$ & 0.40--24.50 & 800 & -87.0 & 14.73 & 227.8 \\
		力学 & $\alpha_{\varepsilon}=1.5$ & 0.40--17.10 & 665 & -138.2 & 17.46 & 256.5 \\
		\midrule
		间隙传热 & $\alpha_h=0.5$ & 0.40--24.65 & 803 & -7.0 & 15.41 & 227.8 \\
		间隙传热 & $\alpha_h=1.0$ & 0.40--24.50 & 800 & -87.0 & 14.73 & 227.8 \\
		间隙传热 & $\alpha_h=2.0$ & 0.40--24.45 & 799 & -155.1 & 16.22 & 230.3 \\
		\midrule
		燃料热导率 & $\alpha_k=0.9$ & 0.40--24.55 & 801 & -26.6 & 13.27 & 215.4 \\
		燃料热导率 & $\alpha_k=1.0$ & 0.40--24.50 & 800 & -87.0 & 14.73 & 227.8 \\
		燃料热导率 & $\alpha_k=1.1$ & 0.40--24.50 & 800 & -134.9 & 16.03 & 232.5 \\
		\bottomrule
	\end{tabular*}
\end{table}

图~\ref{fig:ifa651r1_temperature_alpha_sensitivity} 和表~\ref{tab:ifa651r1_temperature_alpha_sensitivity_error} 给出了 IFA-651-r1 的敏感性分析结果。该燃料棒基准模型整体低估中心温度，MBE 为 \(-87.0\) K。减小燃料几何演化缩放因子可明显改善误差，其中 \(\alpha_{\varepsilon}=0.5\) 时 MBE 降至 \(-18.3\) K，RMSE 也由 227.8 K 降至 192.3 K。减小燃料热导率缩放因子同样能够一定程度上改善低估偏差，而增大 \(\alpha_k\)、\(\alpha_h\) 或 \(\alpha_{\varepsilon}\) 通常会使中心温度进一步降低，导致负偏差扩大。需要注意的是，\(\alpha_{\varepsilon}=1.5\) 时有效对比燃耗范围缩短，说明较强的几何演化扰动可能加剧后期计算的不稳定性。

\begin{figure}[!htb]
	\centering
	
	\begin{subfigure}[t]{0.48\textwidth}
		\centering
		\includegraphics[width=\linewidth]{IFA-651-r3/IFA-651-r3_temperature_fuel_thermal_conductivity_alpha_comparison.pdf}
		\caption{燃料热导率缩放因子}
		\label{fig:ifa651r3_temp_alpha_fuel_thermal_conductivity}
	\end{subfigure}
	\hfill
	\begin{subfigure}[t]{0.48\textwidth}
		\centering
		\includegraphics[width=\linewidth]{IFA-651-r3/IFA-651-r3_temperature_gap_conductivity_alpha_comparison.pdf}
		\caption{间隙传热缩放因子}
		\label{fig:ifa651r3_temp_alpha_gap_conductance}
	\end{subfigure}
	
	\vspace{0.5em}
	
	\begin{subfigure}[t]{0.48\textwidth}
		\centering
		\includegraphics[width=\linewidth]{IFA-651-r3/IFA-651-r3_temperature_geometricstructure_alpha_comparison.pdf}
		\caption{燃料几何演化缩放因子}
		\label{fig:ifa651r3_temp_alpha_geometric_evolution}
	\end{subfigure}
	
	\caption{IFA-651-r3 燃料中心温度对不同物理模型缩放因子的响应}
	\label{fig:ifa651r3_temperature_alpha_sensitivity}
\end{figure}

\begin{table}[H]
	\centering
	\caption{IFA-651-r3 不同缩放因子下中心温度预测误差统计}
	\label{tab:ifa651r3_temperature_alpha_sensitivity_error}
	\small
	\begin{tabular*}{0.9\textwidth}{
			@{\extracolsep{\fill}}
			llccccc
			@{}
		}
		\toprule
		扰动对象 & 缩放因子 & 燃耗范围 & 对比点数 & MBE/K & $\overline{\delta T}$/\% & RMSE/K \\
		\midrule
		力学 & $\alpha_{\varepsilon}=0.5$ & 0.10--23.13 & 457 & -56.7 & 6.15 & 143.3 \\
		力学 & $\alpha_{\varepsilon}=1.0$ & 0.10--23.13 & 457 & -60.0 & 6.70 & 147.4 \\
		力学 & $\alpha_{\varepsilon}=1.5$ & 0.10--23.08 & 456 & -57.5 & 6.96 & 147.4 \\
		\midrule
		间隙传热 & $\alpha_h=0.5$ & 0.10--23.33 & 461 & 86.0 & 11.81 & 167.1 \\
		间隙传热 & $\alpha_h=1.0$ & 0.10--23.13 & 457 & -60.0 & 6.70 & 147.4 \\
		间隙传热 & $\alpha_h=2.0$ & 0.10--22.98 & 454 & -172.1 & 14.52 & 210.5 \\
		\midrule
		燃料热导率 & $\alpha_k=0.9$ & 0.10--23.13 & 457 & -4.9 & 6.69 & 144.8 \\
		燃料热导率 & $\alpha_k=1.0$ & 0.10--23.13 & 457 & -60.0 & 6.70 & 147.4 \\
		燃料热导率 & $\alpha_k=1.1$ & 0.10--23.08 & 456 & -107.0 & 9.32 & 167.3 \\
		\bottomrule
	\end{tabular*}
\end{table}

图~\ref{fig:ifa651r3_temperature_alpha_sensitivity} 和表~\ref{tab:ifa651r3_temperature_alpha_sensitivity_error} 给出了 IFA-651-r3 的敏感性分析结果。该燃料棒基准模型的平均相对误差为 6.70\%，RMSE 为 147.4 K，整体预测精度相对较好。燃料几何演化缩放因子对误差指标影响较小，说明该算例中中心温度对几何演化扰动并不敏感。相比之下，间隙传热缩放因子影响更为明显，降低或提高 \(\alpha_h\) 均会使 MBE 明显偏离零值，并增大平均相对误差和 RMSE。燃料热导率缩放因子中，\(\alpha_k=0.9\) 可使 MBE 接近零，但整体误差指标改善有限。因此，IFA-651-r3 的基准模型已经能够较好描述中心温度演化，不需要较大的参数修正。

\begin{figure}[!htb]
	\centering
	
	\begin{subfigure}[t]{0.48\textwidth}
		\centering
		\includegraphics[width=\linewidth]{IFA-651-r6/IFA-651-r6_temperature_fuel_thermal_conductivity_alpha_comparison.pdf}
		\caption{燃料热导率缩放因子}
		\label{fig:ifa651r6_temp_alpha_fuel_thermal_conductivity}
	\end{subfigure}
	\hfill
	\begin{subfigure}[t]{0.48\textwidth}
		\centering
		\includegraphics[width=\linewidth]{IFA-651-r6/IFA-651-r6_temperature_gap_conductivity_alpha_comparison.pdf}
		\caption{间隙传热缩放因子}
		\label{fig:ifa651r6_temp_alpha_gap_conductance}
	\end{subfigure}
	
	\vspace{0.5em}
	
	\begin{subfigure}[t]{0.48\textwidth}
		\centering
		\includegraphics[width=\linewidth]{IFA-651-r6/IFA-651-r6_temperature_geometristructure_alpha_comparison.pdf}
		\caption{燃料几何演化缩放因子}
		\label{fig:ifa651r6_temp_alpha_geometric_evolution}
	\end{subfigure}
	
	\caption{IFA-651-r6 燃料中心温度对不同物理模型缩放因子的响应}
	\label{fig:ifa651r6_temperature_alpha_sensitivity}
\end{figure}

{
	\small
	
	\begin{xltabular}{\textwidth}{
			@{\extracolsep{\fill}}
			llccccc
			@{}
		}
		
		\caption{IFA-651-r6 不同缩放因子下中心温度预测误差统计}
		\label{tab:ifa651-r6_temperature_alpha_sensitivity_error}
		\\
		
		\toprule
		扰动对象 & 缩放因子 & 燃耗范围 & 对比点数 & MBE/K & $\overline{\delta T}$/\% & RMSE/K \\
		\midrule
		\endfirsthead
		
		\multicolumn{7}{c}{
			\tablename~\thetable\quad
			IFA-651-r6 不同缩放因子下中心温度预测误差统计（续）
		}
		\\
		
		\toprule
		扰动对象 & 缩放因子 & 燃耗范围 & 对比点数 & MBE/K & $\overline{\delta T}$/\% & RMSE/K \\
		\midrule
		\endhead
		
		\midrule
		\multicolumn{7}{r}{续下页}
		\\
		\endfoot
		
		\bottomrule
		\endlastfoot
		
		力学 & $\alpha_{\varepsilon}=0.5$ & 0.09--22.03 & 496 & -67.9 & 16.49 & 244.5 \\
		力学 & $\alpha_{\varepsilon}=1.0$ & 0.09--21.99 & 495 & -249.9 & 21.98 & 299.6 \\
		力学 & $\alpha_{\varepsilon}=1.5$ & 0.09--21.94 & 494 & -455.8 & 39.59 & 471.1 \\
		
		\midrule
		
		间隙传热 & $\alpha_h=0.5$ & 0.09--22.16 & 499 & -104.7 & 11.96 & 206.4 \\
		间隙传热 & $\alpha_h=1.0$ & 0.09--21.99 & 495 & -249.9 & 21.98 & 299.6 \\
		间隙传热 & $\alpha_h=2.0$ & 0.09--21.81 & 491 & -398.6 & 34.98 & 421.0 \\
		
		\midrule
		
		燃料热导率 & $\alpha_k=0.9$ & 0.09--22.03 & 496 & -200.3 & 18.01 & 265.5 \\
		燃料热导率 & $\alpha_k=1.0$ & 0.09--21.99 & 495 & -249.9 & 21.98 & 299.6 \\
		燃料热导率 & $\alpha_k=1.1$ & 0.09--21.94 & 494 & -296.6 & 25.99 & 335.4 \\
		
	\end{xltabular}
	
}

图~\ref{fig:ifa651r6_temperature_alpha_sensitivity} 和表~\ref{tab:ifa651-r6_temperature_alpha_sensitivity_error} 给出了 IFA-651-r6 的敏感性分析结果。该燃料棒基准模型明显低估中心温度，MBE 为 \(-249.9\) K，平均相对误差为 21.98\%，RMSE 为 299.6 K。减小燃料几何演化、间隙传热和燃料热导率缩放因子均能够提高中心温度并改善误差，其中 \(\alpha_h=0.5\) 的改善最为明显，平均相对误差降低至 11.96\%，RMSE 降低至 206.4 K。相反，增大 \(\alpha_{\varepsilon}\)、\(\alpha_h\) 或 \(\alpha_k\) 会使中心温度进一步降低，导致负偏差显著扩大。该结果表明，IFA-651-r6 的偏差主要与间隙传热能力偏强、燃料几何演化导致的间隙闭合过强或燃料热导率偏高有关。

\subsubsection{中心温度偏差来源综合判断}

综合五根燃料棒的敏感性分析结果可以看出，燃料几何演化和间隙传热模型对中心温度预测误差的影响最为直接，其作用主要通过改变燃料-包壳间隙闭合时间、接触压力和等效间隙热阻来体现。燃料热导率缩放因子的影响相对稳定，主要改变燃料内部径向温降。对于基准模型高估中心温度的算例，提高燃料热导率、增强间隙传热或增强燃料几何演化通常会降低中心温度并改善误差；而对于基准模型低估中心温度的算例，上述扰动可能进一步降低中心温度，从而导致误差增大。

因此，缩放因子分析的主要作用不是给出统一的参数修正方向，而是识别不同物理过程对中心温度偏差的敏感性及其燃耗阶段特征。总体来看，燃料-包壳间隙状态、接触传热建立过程以及燃料几何演化模型是影响中心温度预测偏差的重要因素，后续模型改进应重点关注燃料重定位、密实化、辐照肿胀及间隙传热模型之间的耦合关系。

\xsection{裂变气体释放模型的验证与对比}
{Validation and Comparison of Fission Gas Release Models}

前文分别建立了ForMas和SCIANTIX两种复杂程度不同的裂变气体行为模型。ForMas模型采用相对简化的扩散--饱和--释放过程描述，根据燃料温度、裂变率、燃耗和晶粒尺度等状态参数计算裂变气体由晶内向晶界迁移及其后续释放；SCIANTIX则进一步描述气体原子的晶内扩散、再溶解与俘获、晶内和晶界气泡演化以及晶界气体释放等细观过程。两种模型在物理描述层次和计算成本方面存在明显差异，因此有必要在统一的燃料棒热工--力学计算框架下，对其裂变气体释放预测能力进行系统比较。

本节选取FRAPCON-4.0积分评估数据库中的MOX燃料棒辐照实验作为主要验证对象，并按照运行历史将案例划分为长期稳态辐照案例和寿期末功率提升案例。对于具有实验裂变气体释放率（fission gas release, FGR）测量结果的燃料棒，以实验寿期末FGR作为定量验证基准，同时将FRAPCON-4.0计算结果作为参考。对于ForMas和SCIANTIX输出的完整FGR时间历程，则结合燃料功率和温度历史分析模型对运行条件变化的响应特征。在此基础上，进一步讨论不同模型的预测规律、主要偏差来源及其适用范围。


\xsubsection{稳态辐照案例FGR验证与模型对比}
{FGR Validation and Model Comparison under Steady-State Irradiation}

本节所称的“稳态辐照”并非指燃料棒在整个辐照寿期内始终保持恒定功率，而是指燃料棒主要经历正常运行条件下的长时程辐照。其功率和冷却剂边界条件可以随循环运行、停堆及功率调节发生变化，但不包含分钟级快速瞬态或专门设置的寿期末功率提升过程。按照FRAPCON-4.0积分评估报告的分类，其MOX燃料评估数据库同时包括长期稳态功率运行燃料棒和在稳态辐照基础上进一步经历寿期末功率提升的燃料棒，两类案例分别用于评价正常辐照条件和功率快速变化条件下的裂变气体释放行为。

表~\ref{tab:steady_state_mox_fgr_comparison}汇总了本文选取的18组稳态MOX燃料棒案例，包括Halden反应堆IFA系列、美国先进试验堆（ATR）辐照棒以及M504和M308等试验燃料棒。所选案例的棒平均燃耗约为20--68.3~GWd/MTU，实验FGR由不足1\%延伸至约22\%，覆盖了较宽的燃耗和释放率范围。计算过程中采用各燃料棒相应的功率历史和热工边界条件，而非采用单一恒定功率近似整个辐照过程。

FRAPCON-4.0积分评估报告指出，全部稳态MOX案例的预测误差标准差约为6.8个百分点。部分ATR燃料棒由于存在较明显的径向中子通量梯度，棒平均功率的确定具有较大不确定性；排除这些ATR案例后，FRAPCON-4.0稳态MOX裂变气体释放预测的绝对标准差降低至约4.4个百分点。报告同时指出，在棒平均燃耗不超过约62~GWd/MTU时，FRAPCON-4.0能够对较宽范围的MOX燃料FGR给出合理预测，而在更高燃耗区域存在逐渐低估FGR的可能。因此，IFA-629-3r5和IFA-629-3r6除用于评价模型总体预测能力外，还可用于考察不同FGR模型在较高燃耗条件下的外推特征。

对于IFA-651.1系列燃料棒，表中带星号的实验FGR并非由辐照后穿刺实验直接测得，而是根据燃料棒内压数据反演得到，其不确定性通常高于穿刺法直接测量结果。因此，后续模型误差分析同时考虑实验数据来源和运行历史的不确定性，不依据单个燃料棒的预测误差对模型性能作出判断。

\begin{table}[htbp]
	\centering
	\caption{稳态辐照MOX燃料棒FGR实验值与不同模型预测结果}
	\label{tab:steady_state_mox_fgr_comparison}
	\small
	\begin{tabularx}{\textwidth}{
			>{\raggedright\arraybackslash}X c c c c c}
		\toprule
		Rod &
		\begin{tabular}[c]{@{}c@{}}
			平均燃耗\\
			GWd/MTU
		\end{tabular} &
		\begin{tabular}[c]{@{}c@{}}
			实验FGR\\
			\%
		\end{tabular} &
		\begin{tabular}[c]{@{}c@{}}
			FRAPCON-4.0\\
			\%
		\end{tabular} &
		\begin{tabular}[c]{@{}c@{}}
			ForMas\\
			\%
		\end{tabular} &
		\begin{tabular}[c]{@{}c@{}}
			SCIANTIX\\
			\%
		\end{tabular} \\
		\midrule
		
		IFA-651.1r1               & 22.41 & $10^{*}$ & 1.13  & 4.73  & 8.90  \\
		IFA-651.1r3               & 21.73 & $2^{*}$  & 1.68  & 4.31  & 8.85  \\
		IFA-651.1r6               & 20.27 & $7^{*}$  & 1.88  & 2.61  & 4.22  \\
		ATR PII C2 P5             & 21.00 & 1.146    & 22.19 & 6.25  & 14.20 \\
		ATR PIII C3 P6            & 30.00 & 1.253    & 15.96 & 4.36  & 11.14 \\
		ATR PIII C10 P13          & 30.00 & 2.019    & 16.98 & 4.74  & 12.20 \\
		ATR PIV C4 P7             & 40.00 & 8.214    & 18.69 & 4.37  & 5.03  \\
		ATR PIV C5 P8             & 50.00 & 3.009    & 7.12  & 1.19  & 4.23  \\
		ATR PIV C6 P9             & 50.00 & 7.066    & 10.11 & 2.07  & 4.82  \\
		ATR PIV C12 P15           & 50.00 & 8.761    & 9.87  & 1.95  & 5.39  \\
		IFA-629.1                 & 29.00 & 21.7     & 19.13 & 5.77  & 18.41 \\
		M504 H8                   & 37.50 & 0.54     & 0.26  & 1.24  & 1.190 \\
		M504 I2                   & 43.00 & 0.85     & 0.75  & 1.77  & 1.77  \\
		IFA-597.4/.5/.6/.7r10     & 35.70 & 17.0     & 13.17 & 4.65  & 15.29 \\
		IFA-597.4/.5/.6/.7r11     & 36.80 & 14.0     & 20.95 & 3.54  & 10.28 \\
		IFA-629-3r5               & 68.30 & 21.0     & 6.28  & 6.74  & 16.96 \\
		IFA-629-3r6               & 63.60 & 12.0     & 5.73  & 36.24 & 24.58 \\
		M308 Segment 2            & 57.50 & 5.0      & 4.24  & 0.10  & 0.10  \\
		
		\bottomrule
		\multicolumn{6}{l}{
			\footnotesize
			$^{*}$基于棒内压力数据推算的寿期末FGR，
			其不确定性高于穿刺法直接测量结果。}
	\end{tabularx}
\end{table}

考虑到部分实验FGR低于1\%，采用传统相对误差会显著放大小释放率案例的误差权重，因此本文主要采用FGR绝对偏差及以百分点为单位的统计指标评价模型预测性能。对于第$i$个燃料棒，定义预测偏差为
\begin{equation}
	\Delta F_i
	=
	F_{i,\mathrm{cal}}-F_{i,\mathrm{exp}},
	\label{eq:fgr_deviation}
\end{equation}
式中，$F_{i,\mathrm{cal}}$和$F_{i,\mathrm{exp}}$分别为计算和实验测得的FGR，单位均为\%。进一步定义平均偏差（MBE）、平均绝对偏差（MAE）和均方根误差（RMSE）分别为
\begin{equation}
	\mathrm{MBE}
	=
	\frac{1}{N}
	\sum_{i=1}^{N}\Delta F_i,
	\label{eq:fgr_mbe}
\end{equation}
\begin{equation}
	\mathrm{MAE}
	=
	\frac{1}{N}
	\sum_{i=1}^{N}
	\left|\Delta F_i\right|,
	\label{eq:fgr_mae}
\end{equation}
\begin{equation}
	\mathrm{RMSE}
	=
	\sqrt{
		\frac{1}{N}
		\sum_{i=1}^{N}
		\Delta F_i^2},
	\label{eq:fgr_rmse}
\end{equation}
其中$N$为参与统计的燃料棒数量。

根据表~\ref{tab:steady_state_mox_fgr_comparison}中的18组案例，三种模型的误差统计结果见表~\ref{tab:steady_state_mox_fgr_statistics}。

\begin{table}[htbp]
	\centering
	\caption{稳态辐照MOX燃料棒FGR预测误差统计}
	\label{tab:steady_state_mox_fgr_statistics}
	\begin{tabularx}{\textwidth}{
			>{\raggedright\arraybackslash}Xccccc}
		\toprule
		模型 &
		MBE/百分点 &
		MAE/百分点 &
		RMSE/百分点 &
		$|\Delta F|\leq5$ &
		$|\Delta F|\leq10$ \\
		\midrule
		FRAPCON-4.0 & 1.86  & 6.62 & 8.98 & 9/18  & 13/18 \\
		ForMas      & -2.55 & 6.90 & 9.18 & 10/18 & 13/18 \\
		SCIANTIX    & 1.39  & 4.76 & 6.16 & 13/18 & 15/18 \\
		\bottomrule
	\end{tabularx}
\end{table}

由表~\ref{tab:steady_state_mox_fgr_statistics}可见，ForMas模型在全部18组稳态案例中的MBE为$-2.55$个百分点，表现出一定的整体负偏差，其MAE和RMSE分别为6.90和9.18个百分点，与FRAPCON-4.0的总体误差量级接近。18组案例中有10组ForMas预测结果与实验值的绝对偏差不超过5个百分点，13组不超过10个百分点，表明该简化模型能够在较低计算复杂度下给出多数案例的FGR数量级。

然而，ForMas在不同释放水平下的偏差分布并不均匀。对于IFA-629.1、IFA-597系列和IFA-629-3r5等实验FGR较高的案例，ForMas均表现出不同程度的低估。例如IFA-629.1实验FGR为21.7\%，ForMas预测值仅为5.77\%；IFA-597.4/.5/.6/.7r10实验值为17.0\%，预测值为4.65\%。这一结果表明，ForMas模型虽然能够描述正常辐照条件下裂变气体逐步释放的总体趋势，但对于部分高释放状态，其气体向晶界迁移和后续快速释放过程的响应仍然偏弱。

与ForMas相比，SCIANTIX在该组稳态案例中的总体预测性能有所改善，其MBE、MAE和RMSE分别为1.39、4.76和6.16个百分点，其中13组案例的绝对偏差不超过5个百分点，15组不超过10个百分点。改善主要体现在较高释放率案例。例如IFA-629.1的预测值由ForMas的5.77\%提高至18.41\%，IFA-597.4/.5/.6/.7r10由4.65\%提高至15.29\%，IFA-629-3r5则由6.74\%提高至16.96\%，均明显向实验结果靠近。这表明在较高裂变气体释放水平下，对晶内气体迁移、气泡演化及晶界气体行为进行更细致描述有助于提高FGR预测能力。

需要指出的是，模型复杂程度增加并不意味着所有案例的误差均单调减小。部分低释放ATR案例中，SCIANTIX仍存在明显高估；IFA-629-3r6则表现出与大部分高释放案例不同的异常特征，其实验FGR为12.0\%，ForMas和SCIANTIX分别预测为36.24\%和24.58\%。该案例说明FGR预测不仅受燃耗水平控制，还对具体功率历史、燃料温度以及模型内部释放判据具有较强敏感性。对于这类异常案例，需要结合完整的功率--温度--FGR时间历程进一步分析，而不能仅依据寿期末燃耗判断其偏差来源。


\xsubsection{功率提升工况FGR终值验证与动态响应}
{End-of-Life Validation and Dynamic FGR Response under Power-Ramp Conditions}

为进一步评价裂变气体模型对快速功率变化工况的响应能力，选取M501系列8根MOX燃料棒进行计算。与前述稳态案例不同，M501燃料棒在长期辐照后经历了寿期末短时间高功率运行。其中HR系列和MR-1至MR-3的高功率阶段持续约12~h，而MR-4的高功率阶段仅持续约20~min，因此该组案例能够同时考察功率提升幅度和高温持续时间对裂变气体释放的影响。

需要说明的是，M501实验仅提供功率提升结束后的总FGR测量结果，而未给出与计算相对应的完整FGR时间历程。因此，本节以寿期末实验FGR作为模型的直接定量验证依据；ForMas和SCIANTIX计算获得的完整FGR时间历程则用于分析两种模型自身对功率和温度变化的动态响应特征，而不将计算过程曲线视为经过实验直接验证的瞬态FGR历史。

M501案例的实验测量值、FRAPCON-4.0预测值及ForMas计算结果见表~\ref{tab:m501_fgr_comparison}。SCIANTIX对应结果将在采用相同功率历史和边界条件完成计算后纳入统一比较。

\begin{table}[htbp]
	\centering
	\caption{M501功率提升MOX燃料棒FGR实验值与不同模型预测结果}
	\label{tab:m501_fgr_comparison}
	\small
	\begin{tabularx}{\textwidth}{
			>{\raggedright\arraybackslash}X c c c c c c c}
		\toprule
		燃料棒 &
		\begin{tabular}[c]{@{}c@{}}
			平均燃耗\\
			GWd/MTU
		\end{tabular} &
		\begin{tabular}[c]{@{}c@{}}
			最高线功率\\
			kW$\cdot$m$^{-1}$
		\end{tabular} &
		\begin{tabular}[c]{@{}c@{}}
			高功率\\
			持续时间
		\end{tabular} &
		实验FGR/\% &
		FRAPCON-4.0/\% &
		ForMas/\% &
		SCIANTIX/\% \\
		\midrule
		
		M501 HR-1 & 37 & 38.1 & 12 h   & 7.67  & 19.55 & 3.901 & 2.725 \\
		M501 HR-2 & 37 & 35.7 & 12 h   & 8.24  & 16.92 & 4.758 & 7.701 \\
		M501 HR-3 & 37 & 46.2 & 12 h   & 18.21 & 29.87 & 4.339 & 5.720 \\
		M501 HR-4 & 36 & 47.0 & 12 h   & 16.04 & 31.63 & 4.547 & -- \\
		
		M501 MR-1 & 34 & 38.1 & 12 h   & 2.43  & 12.32 & 3.410 & -- \\
		M501 MR-2 & 34 & 41.9 & 12 h   & 9.20  & 19.40 & 3.604 & -- \\
		M501 MR-3 & 34 & 40.5 & 12 h   & 6.39  & 18.43 & 3.164 & -- \\
		M501 MR-4 & 33 & 41.7 & 20 min & 2.17  & 5.00  & 3.383 & -- \\
		
		\bottomrule
	\end{tabularx}
\end{table}

在目前已获得的ForMas结果中，8组M501案例的MBE、MAE和RMSE分别为$-4.91$、5.45和7.03个百分点，而FRAPCON-4.0对应结果分别为10.35、10.35和10.90个百分点。虽然ForMas在该组案例中的总体绝对误差低于FRAPCON-4.0，但其预测结果主要集中在3.16\%--4.76\%范围内，而实验FGR实际覆盖2.17\%--18.21\%。因此，较低的总体误差并不能说明ForMas已经准确描述了功率提升过程中的裂变气体释放机制。

特别是M501 HR-3和HR-4，其实验FGR分别达到18.21\%和16.04\%，而ForMas仅预测为4.339\%和4.547\%，表明该模型对高功率提升条件下附加气体释放的响应明显不足。相反，对于实验释放量较低的MR-1和MR-4，ForMas预测值与实验结果的差异相对较小。这种预测值向较窄范围聚集的特征说明，仅依据寿期末总FGR还不足以判断模型偏差究竟来源于前期气体积累还是功率提升阶段的附加释放，因此有必要利用完整时间历程进一步分解。

为定量描述功率提升阶段对寿期末FGR的贡献，定义功率提升开始前的FGR为$F_{\mathrm{pre}}$，寿期末FGR为$F_{\mathrm{EOL}}$，则功率提升阶段新增释放量定义为
\begin{equation}
	\Delta F_{\mathrm{ramp}}
	=
	F_{\mathrm{EOL}}-F_{\mathrm{pre}},
	\label{eq:fgr_ramp_increment}
\end{equation}
相应的功率提升贡献比例定义为
\begin{equation}
	\eta_{\mathrm{ramp}}
	=
	\frac{\Delta F_{\mathrm{ramp}}}
	{F_{\mathrm{EOL}}}
	\times100\%.
	\label{eq:fgr_ramp_fraction}
\end{equation}

式~\eqref{eq:fgr_ramp_increment}用于反映功率提升阶段新增的裂变气体释放量，而式~\eqref{eq:fgr_ramp_fraction}用于衡量该阶段对最终FGR的相对贡献。完成SCIANTIX功率提升案例计算后，可分别比较ForMas和SCIANTIX的$F_{\mathrm{pre}}$、$\Delta F_{\mathrm{ramp}}$和$\eta_{\mathrm{ramp}}$，从而判断两种模型在长期气体积累和短时高温附加释放两个阶段的差异。

为了进一步分析模型动态响应，可选取具有代表性的M501 MR-1、HR-1、HR-3和MR-4案例，对其功率、燃料中心温度和FGR时间历程进行对比。其中MR-1代表低释放水平，HR-1代表中等释放水平，HR-3代表高释放水平，而MR-4由于高功率阶段仅持续20~min，可用于分析高温持续时间对裂变气体释放的影响。

% SCIANTIX功率提升案例计算完成后，可加入如下图：
%
% \begin{figure}[htbp]
	% 	\centering
	% 	\includegraphics[width=0.95\textwidth]
	% 	{Figures/m501_fgr_dynamic_response.pdf}
	% 	\caption{典型M501功率提升案例的功率、燃料温度及FGR动态响应}
	% 	\label{fig:m501_fgr_dynamic_response}
	% \end{figure}
%
% 图中建议同时给出：
% 1. 线功率历史；
% 2. 燃料中心温度历史；
% 3. ForMas和SCIANTIX的FGR时间历程；
% 4. 寿期末实验FGR以散点形式标出。


\xsubsection{不同运行条件下的FGR预测规律}
{FGR Prediction Characteristics under Different Operating Conditions}

上述稳态和功率提升案例表明，裂变气体释放并不能由寿期末燃耗这一单一变量描述，而是燃耗累积、燃料温度水平、功率历史以及高温持续时间共同作用的结果。为避免仅依据个别案例判断模型优劣，本节从释放水平、燃耗和功率变化三个方面总结两种裂变气体模型的预测特征。

首先，从实验FGR水平与预测偏差的关系看，ForMas在中低释放率案例中能够给出与实验结果相同的数量级，但随着实验FGR升高，其低估现象总体更加明显。IFA-629.1、IFA-597系列、IFA-629-3r5以及M501 HR-3和HR-4均具有这一特征。相比之下，SCIANTIX在多组较高释放率稳态案例中明显减小了预测偏差，说明更加细致的晶内--晶界气体迁移与气泡演化描述增强了模型对高释放状态的响应能力。但IFA-629-3r6等异常案例同时说明，预测误差并非实验FGR的简单单调函数，具体运行历史仍具有重要影响。

% 建议加入实验值-预测值散点图
%
% \begin{figure}[htbp]
	% 	\centering
	% 	\includegraphics[width=0.72\textwidth]
	% 	{Figures/fgr_prediction_vs_measurement.pdf}
	% 	\caption{不同裂变气体模型预测FGR与实验结果对比}
	% 	\label{fig:fgr_prediction_vs_measurement}
	% \end{figure}
%
% 横坐标：实验FGR
% 纵坐标：计算FGR
% 同时绘制y=x参考线
% 数据系列：FRAPCON-4.0、ForMas、SCIANTIX

其次，从燃耗与预测误差的关系看，当前案例未表现出“燃耗越高、FGR越高”或“燃耗越高、预测误差越大”的简单单调关系。在约30--40~GWd/MTU的相近燃耗区间内，实验FGR既包含低于1\%的M504案例，也包含超过10\%甚至20\%的IFA案例；类似地，在较高燃耗区域，不同燃料棒的释放水平仍存在明显差异。这说明燃耗主要反映裂变气体的累积生成程度，而实际释放量还取决于燃料在整个寿期内所经历的温度历史。由于裂变气体扩散和气泡演化均具有明显的温度依赖性，相同或接近燃耗条件下，不同功率历史仍可能产生显著不同的FGR。

% 建议加入燃耗-偏差关系图
%
% \begin{figure}[htbp]
	% 	\centering
	% 	\includegraphics[width=0.72\textwidth]
	% 	{Figures/fgr_error_vs_burnup.pdf}
	% 	\caption{不同模型FGR预测偏差随燃耗的变化}
	% 	\label{fig:fgr_error_vs_burnup}
	% \end{figure}

再次，对于M501功率提升案例，寿期末FGR不仅取决于功率提升前已经积累并释放的气体，还取决于高功率阶段产生的附加释放。功率提高使燃料温度升高，从而增强裂变气体向晶界的迁移，并可能使已有晶界气体在较短时间内进一步释放。因此，对功率提升工况的评价不能仅比较最终FGR，还应同时考察$\Delta F_{\mathrm{ramp}}$及$\eta_{\mathrm{ramp}}$。

M501 MR-4与其余案例的高功率持续时间差异尤其具有分析意义。其高功率阶段仅持续20~min，而其余案例主要持续12~h。对于扩散控制的裂变气体释放过程，即使短时间内达到较高燃料温度，有限的高温持续时间仍会限制气体由晶内向晶界迁移的累积过程。因此，在比较不同功率提升案例时，应同时考虑峰值功率、对应的燃料峰值温度以及高温持续时间，而不能仅依据最高线功率判断其最终释放程度。

综合上述结果，本文认为裂变气体释放的主要控制关系可概括为
\begin{equation}
	FGR
	=
	f\left[
	B(t),\,
	T(t),\,
	\dot{f}(t),\,
	d_{\mathrm{g}},\,
	t_{\mathrm{high}}
	\right],
	\label{eq:fgr_operating_dependence}
\end{equation}
式中，$B(t)$为燃耗历史，$T(t)$为燃料温度历史，$\dot{f}(t)$为裂变率历史，$d_{\mathrm{g}}$为晶粒尺度，$t_{\mathrm{high}}$为高温持续时间。式~\eqref{eq:fgr_operating_dependence}并非新的经验关联式，而是用于概括本文验证结果所体现的主要影响因素。其中，燃耗决定裂变气体的累积生成水平，而温度和时间历史则直接影响气体迁移与释放过程。


\xsubsection{模型差异、偏差来源与适用范围}
{Model Differences, Deviation Sources, and Applicability}

从上述稳态和功率提升案例可以看出，ForMas与SCIANTIX的主要差异并不仅体现为寿期末FGR数值的不同，更来源于两种模型对裂变气体迁移和释放过程描述层次的差异。

ForMas采用简化的扩散--饱和--释放框架，其优势在于所需状态变量较少、计算开销较低，适合在燃料棒全寿期多物理场分析中快速预测裂变气体释放及其对内压和间隙传热的反馈。稳态案例结果表明，ForMas能够在多数中低释放水平工况下给出合理数量级的预测，其总体MAE与FRAPCON-4.0处于相近水平。然而，在IFA-629.1、IFA-597系列和M501高释放案例中存在明显低估，说明其对于部分高温、高释放状态的响应仍然不足。

这种偏差可能首先与有效扩散过程的简化有关。裂变气体由晶内向晶界迁移对温度具有较强敏感性，当模型使用等效扩散系数描述多种微观迁移机制时，扩散参数的选取将直接影响气体到达晶界的速率。对于快速升温或较高温度条件，若等效扩散过程不能充分反映实际气体迁移增强，则会导致晶界气体库存和最终释放量偏低。

其次，晶粒尺寸直接影响裂变气体到达晶界所需的特征扩散距离。对于不同燃料棒，如果采用统一或理想化的初始晶粒尺寸，而未充分考虑制造状态及辐照过程中的晶粒演化，则可能导致气体迁移时间尺度出现系统偏差。因此，晶粒尺寸不仅影响寿期末FGR，对功率快速变化阶段的释放响应同样具有重要作用。

再次，ForMas中的晶界饱和和释放判据决定了已经迁移至晶界的气体何时进入自由空间。对于长期稳态辐照，该简化处理可以较有效地描述FGR的逐步增长；但在短时间高功率运行条件下，实际晶界气泡连通、长大及破裂过程可能表现出更强的非线性。如果释放判据对温度或晶界库存变化的响应不足，则容易形成M501案例中计算FGR集中于较窄范围的现象。

SCIANTIX通过显式描述晶内气体、晶内气泡、晶界气体及相关演化过程，在多数高释放稳态案例中改善了预测结果，其总体MAE由ForMas的6.90个百分点降低至4.76个百分点。然而，SCIANTIX仍未消除全部预测偏差，部分低释放ATR案例仍表现出高估，IFA-629-3r6也存在较明显偏高现象。这表明增加微观模型复杂程度能够改善部分气体释放机制的描述，但模型结果仍受到温度场、裂变率、初始晶粒尺寸、模型参数以及实验功率历史不确定性的共同影响。

对于高燃耗案例，还需要考虑长期辐照导致的燃料微观结构变化。随着燃耗提高，裂变产物积累、气泡演化以及局部微结构变化可能改变裂变气体扩散和滞留行为。ForMas对这些过程采用较强的等效化处理，因此其高燃耗预测更依赖于经验参数和释放准则；SCIANTIX虽然包含更多微观状态变量，但仍需要合理的材料参数和初始条件才能充分发挥其模型优势。

因此，两种模型的适用性不宜简单概括为“SCIANTIX优于ForMas”。ForMas具有较低的计算成本和较好的多物理场耦合效率，适用于大规模燃料棒全寿期计算、参数研究以及对FGR总体趋势的快速预测；当研究重点涉及较高裂变气体释放、高温历史或快速功率变化条件下的气体迁移与释放过程时，SCIANTIX更加完整的微观机制描述具有明显优势。对于超出主要验证数据库范围的高燃耗或特殊瞬态工况，两种模型的结果均应结合温度历史、功率历史和相关实验不确定性进行综合评价。

完成M501系列SCIANTIX功率提升计算后，可进一步利用两种模型的$F_{\mathrm{pre}}$、$\Delta F_{\mathrm{ramp}}$和$\eta_{\mathrm{ramp}}$进行对比。如果SCIANTIX能够在HR-3和HR-4等高释放案例中给出更大的功率提升阶段附加释放，同时寿期末FGR更加接近实验结果，则可进一步说明ForMas的主要限制来自快速高温阶段的释放响应；反之，如果两种模型均无法再现实验中的高释放水平，则需要进一步从温度边界、裂变率历史、晶粒尺寸以及实验工况不确定性等方面分析偏差来源。
